\chapter{Introduction}

% Enable paragraph indentation
\setlength{\parindent}{1em}

\normalsize

Money flows through modern life in dozens of small, fast events: a coffee, a
taxi, a grocery run, a transfer from a family member, a monthly salary. Each of
these is trivial on its own, yet collectively they determine whether a person
understands and controls their finances. The premise of every personal-finance
application is that if these events are \emph{recorded}, they can be summarized,
categorized, and reflected upon. The recurring failure of those applications is
equally simple: people will not record what is tedious to record. This thesis
presents \textbf{MiPay}, a bilingual, voice-first personal-finance tracker
designed around a single hypothesis---that the dominant obstacle to consistent
expense tracking is \emph{data-entry friction}, and that conversational speech,
processed by a self-hosted artificial-intelligence pipeline, can remove most of
that friction while preserving the user's privacy and incurring no per-request
cost.

This chapter establishes the problem the project addresses, the background and
motivation that make it timely, the enabling technologies, the spectrum of
technical approaches that were considered, the challenges inherent to the task,
and finally the contributions and organization of the thesis.

\section{Problem Statement}

A personal-finance application is only as useful as the data the user is willing
to feed it. Manual transaction entry---selecting a category from a dropdown,
typing an amount, choosing a date, naming a merchant---imposes a cognitive and
physical cost at exactly the wrong moment: while standing at a till, sitting in a
taxi, or walking out of a shop. The result is well documented in product
analytics across the category: high installation rates followed by steep
abandonment, because the marginal effort of logging each transaction exceeds the
perceived marginal benefit.

The problem is sharper still for \textbf{bilingual users in the Arab world}, the
target population of this project. A user in the Gulf or the Levant routinely
mixes Arabic and English within a single spoken sentence---``\,dafact 50
riyal cala \emph{groceries} from Carrefour \emph{yesterday}\,''---a phenomenon
known as \emph{code-switching}~\cite{sitaram2019codeswitch}. Arabic additionally
presents right-to-left script, a rich system of dialects that diverge
substantially from Modern Standard Arabic, and its own set of Arabic-Indic
numerals. A finance application that is to feel natural to this population must
therefore accept free-form bilingual speech, not constrained menu input.

Concretely, the problem this thesis solves can be stated as follows:

\begin{quote}
\itshape
Given a short, spontaneous, possibly code-switched Arabic/English voice
utterance describing one or more financial events, automatically produce a
structured, validated, and editable set of transaction records---each carrying a
type, amount, currency, category, counterparty, and resolved calendar
date---using only self-hosted, open-source models, with latency low enough for
interactive use and with a human confirmation step that prevents any incorrect
record from being silently persisted.
\end{quote}

\section{Background and Motivation}

\subsection{The Friction Problem in Personal Finance}

The behavioral economics of expense tracking are unforgiving. The benefit of a
logged transaction is \emph{deferred and diffuse} (a better monthly summary,
some weeks later), while the cost is \emph{immediate and concrete} (fifteen
seconds of fiddling with a form now). Whenever cost is front-loaded and benefit
is back-loaded, consistent behavior collapses. Reducing the per-event cost from
fifteen seconds of structured input to two seconds of speaking therefore does
not merely make the application more pleasant; it changes whether the core
behavior happens at all. Global financial-inclusion data underline why lowering
this barrier matters: account ownership and digital-payment adoption have risen
sharply across developing economies, yet the tools for everyday money management
remain poorly localized for Arabic speakers~\cite{worldbank2021findex}.

\begin{figure}[H]
    \centering
    \begin{tikzpicture}[
        field/.style={draw, rounded corners, fill=gray!10, minimum width=3.6cm, minimum height=0.8cm, font=\small, align=center},
        stepcircle/.style={circle, draw, fill=black, text=white, minimum size=0.55cm, font=\small\bfseries},
        panellabel/.style={font=\bfseries\large},
        micbody/.style={rounded corners=6pt, draw, fill=black, minimum width=1.0cm, minimum height=1.8cm},
    ]

    % Panel titles
    \node[panellabel] at (0,6) {Manual Entry};
    \node[panellabel] at (8,6) {Voice Input};

    % ---- LEFT PANEL: 5 stacked form fields, each a separate step ----
    \foreach \i/\fieldname in {1/Amount, 2/Category, 3/Currency, 4/Date, 5/Merchant} {
      \pgfmathsetmacro{\y}{5.2 - (\i - 1) * 1.05}
      \node[stepcircle] at (-1.9,\y) {\i};
      \node[field] at (0.2,\y) {\fieldname};
    }
    \node[font=\small\itshape] at (0.2,-0.3) {5 separate steps, roughly 15 seconds};

    % ---- Divider ----
    \draw[dashed, gray, thick] (4,6.5) -- (4,-0.8);

    % ---- RIGHT PANEL: one spoken sentence -> mic -> ready transaction ----
    \node[draw, rounded corners, fill=blue!6, text width=4.6cm, font=\small\itshape,
          align=center, minimum height=1.2cm] (bubble) at (8,5.0)
      {``I paid 50 riyals for groceries at Carrefour yesterday''};

    \node[stepcircle] at (5.9,3.2) {1};
    \node[micbody] (mic) at (8,3.2) {};
    \node[text=white, font=\bfseries\small] at (8,3.2) {MIC};

    \draw[-{Latex[length=2mm]}, thick] (bubble.south) -- (mic.north);

    \node[draw, rounded corners, fill=green!10, minimum width=3.6cm,
          minimum height=0.8cm, font=\small, align=center] (result) at (8,1.4)
      {Transaction ready to confirm};
    \draw[-{Latex[length=2mm]}, thick] (mic.south) -- (result.north);

    \node[font=\small\itshape] at (8,-0.3) {1 step, roughly 2 seconds};

    \end{tikzpicture}
    \caption[Data-entry friction: manual form versus voice input]{Data-entry
    friction compared. Manual entry requires the user to traverse five separate
    interactive fields (amount, category, currency, date, merchant) for every
    transaction, whereas voice input collapses the interaction to a single tap
    and a spoken sentence, deferring all structuring to the AI pipeline.}
    \label{fig:friction}
\end{figure}

\subsection{Why Voice, and Why Now}

Two independent trends make a voice-first design viable in a way it was not a few
years ago. First, \textbf{automatic speech recognition (ASR)} has crossed a
quality threshold for multilingual, accented, and noisy input. The Whisper family
of models, trained on 680{,}000 hours of weakly supervised multilingual
audio~\cite{radford2023whisper}, transcribes Arabic and English---and mixtures of
the two---at a level that was previously the exclusive domain of large commercial
services. Second, \textbf{open-weight large language models (LLMs)} have made it
possible to perform robust natural-language understanding \emph{locally}. A
7-billion-parameter instruction-tuned model such as
Qwen2.5-7B-Instruct~\cite{qwen2024qwen25} can follow detailed extraction
instructions across languages while running on a commodity workstation. The
conjunction of these two trends is what enables a privacy-preserving,
zero-marginal-cost, self-hosted voice pipeline; neither alone would suffice.

Voice assistants themselves are not new---Siri, Alexa, and Google Assistant
popularized the modality~\cite{hoy2018voice,kepuska2018assistants}---but they are
overwhelmingly cloud-dependent and general-purpose. MiPay differs in being
\emph{domain-specific}, \emph{bilingual by design}, and \emph{self-hosted}, which
turns the modality from a convenience into a defensible engineering and research
contribution.

\subsection{The Bilingual and Code-Switching Challenge}

Arabic is the native language of roughly 400 million people and is, from a
computational standpoint, a high-difficulty language: it is morphologically rich,
diglossic (a formal written standard coexisting with many spoken dialects), and
written in a cursive right-to-left script with optional
diacritics~\cite{antoun2020arabert,obeid2020camel}. Spoken financial language in
the region is rarely ``pure'' anything: amounts may be spoken in Arabic while the
merchant name is English; the currency word \emph{riyal} maps to different ISO
codes depending on national context; and Arabic-Indic digits appear alongside
Western digits. A system that treats Arabic and English as two separate,
switchable pipelines will fail on the most natural utterances. MiPay instead
treats code-switched bilingual speech as the \emph{default} case and designs
every stage of the pipeline---transcription, extraction, and
post-processing---to tolerate it.

\subsection{Gap in Existing Solutions}

Existing solutions fall into three groups, each with a structural gap.
\textbf{Commercial budgeting apps} (e.g., the well-known Western expense
trackers) are built around manual or bank-feed input and offer little or no
Arabic voice support; bank-feed integration is moreover largely unavailable in
much of the target region. \textbf{General voice assistants} can transcribe
speech but do not produce structured, persisted financial records, and they ship
the user's audio to remote servers. \textbf{Cloud AI ``wrappers''}---thin
applications that forward audio and text to paid speech and language APIs---can
in principle produce structured output, but they incur per-request monetary cost,
surrender user financial data to third parties, and, from an academic standpoint,
contribute no model-selection or evaluation methodology of their own. The gap
MiPay fills is the intersection: \emph{structured financial extraction, from
bilingual voice, entirely on self-hosted open models, with measurable quality.}

\subsection{Why This Project Matters}

Beyond the immediate product, the project matters for three reasons. From a
\textbf{user perspective}, it lowers the behavioral barrier to financial
self-awareness for an under-served bilingual population. From an
\textbf{engineering perspective}, it demonstrates a complete,
production-shaped, privacy-preserving AI system---speech, language, persistence,
and mobile UI---assembled entirely from open components and reproducible with a
single \texttt{docker compose up}. From a \textbf{research perspective}, the
self-hosted choice is what makes the work a thesis rather than an integration
exercise: it forces explicit model selection, exposes measurable metrics (Word
Error Rate, field-level F\textsubscript{1}), and admits principled ablation
studies (Section~\ref{sec:ch1-contributions}).

\section{Project Aim and Objectives}

The \textbf{aim} of the project is to design, implement, and evaluate a
bilingual, voice-first personal-finance tracker whose AI pipeline is entirely
self-hosted and open-source. This aim decomposes into the following
\textbf{objectives}:

\begin{enumerate}[label=\textbf{O\arabic*.}, leftmargin=*]
    \item Design an end-to-end pipeline that converts a bilingual voice utterance
    into one or more structured, validated transaction records.
    \item Select and justify an open-source ASR model and an open-weight LLM that
    meet quality targets on commodity hardware without a GPU.
    \item Guarantee machine-parseable extraction output through JSON-schema
    constrained decoding, eliminating format-level failure.
    \item Implement a deterministic, unit-tested post-processing layer for the
    sub-tasks at which LLMs are unreliable (numeral normalization, relative-date
    resolution, currency defaulting, and confidence-based review routing).
    \item Build a cross-platform mobile client with full Arabic (right-to-left)
    and English (left-to-right) localization and a human-in-the-loop confirmation
    step.
    \item Construct a labeled bilingual evaluation corpus and a reproducible
    evaluation harness, and report transcription and extraction quality with
    ablation studies.
\end{enumerate}

\section{Scope and Limitations}

The minimum viable product targets the core loop---record, transcribe, extract,
confirm, save---together with manual transaction management, a bilingual
interface, transaction listing with filters, and a simple monthly dashboard. The
following are deliberately \emph{out of scope} for this work and are revisited as
future directions in Chapter~6: budgets, savings goals, recurring-transaction
detection, push notifications, receipt OCR, bank integrations, offline-first
synchronization, multi-currency conversion, and social features. Limiting scope
in this way keeps the thesis focused on its academic core---the self-hosted AI
pipeline and its evaluation---rather than on breadth of product features.

\section{Overview of the Enabling Technologies}

\subsection{Automatic Speech Recognition}

Modern ASR has moved from hybrid hidden-Markov-model/deep-neural-network
systems~\cite{hinton2012deep,povey2011kaldi} through self-supervised acoustic
representation learning~\cite{baevski2020wav2vec} to large, weakly supervised
sequence-to-sequence transformers. Whisper~\cite{radford2023whisper} belongs to
the last category: an encoder--decoder transformer~\cite{vaswani2017attention}
trained to map log-Mel spectrograms directly to text tokens across 99 languages,
with robustness to accent and background noise that emerges from the sheer scale
and diversity of its training data. MiPay uses Whisper through
\texttt{faster-whisper}~\cite{fasterwhisper}, which runs the same weights on the
CTranslate2 inference engine~\cite{ctranslate2} with 8-bit
quantization~\cite{dettmers2022llmint8}, achieving roughly a four-fold speedup on
CPU at negligible accuracy cost. The technology is examined in detail in
Chapter~4.

\subsection{Large Language Models for Information Extraction}

The task of turning ``I paid fifty riyals for groceries at Carrefour yesterday''
into a typed record is an \emph{information-extraction} problem. Historically this
was approached with hand-written grammars or with supervised sequence-labeling
models such as BiLSTM-CRF architectures~\cite{lample2016ner,li2022ner}, both of
which require either brittle rules or large labeled corpora. Instruction-tuned
LLMs~\cite{wei2022flan,ouyang2022instructgpt} changed the economics: a single
well-designed prompt with a handful of in-context
examples~\cite{brown2020gpt3} can perform the extraction across languages and
dialects with no task-specific training data. MiPay adopts this generative,
few-shot approach using a locally hosted open-weight
model~\cite{qwen2024qwen25}, and rigorously bounds its failure modes with the
techniques described next.

\subsection{Constrained Decoding}

A generative model asked to ``answer in JSON'' will, some fraction of the time,
emit malformed JSON, hallucinate extra keys, or choose values outside the allowed
set. \emph{Constrained} (or \emph{guided}) decoding eliminates this by masking,
at every generation step, any token that would violate a supplied grammar or JSON
schema~\cite{willard2023outlines,geng2023grammar}. The output is then
\emph{guaranteed} to be syntactically valid and schema-conformant. MiPay relies
on this guarantee: the LLM is constrained to emit exactly the project's
transaction schema, so the application never needs to defend against unparseable
model output, and enum fields such as \texttt{currency} and \texttt{category} can
only take legal values.

\section{Approaches to Spoken-Language Understanding}

Several families of methods could, in principle, perform the extraction stage.
This section surveys them and explains the hybrid design MiPay adopts.

\subsection{Rule-Based and Grammar Methods}

The earliest and simplest approach is to write explicit rules: regular
expressions for amounts, keyword lists for categories, and pattern grammars for
dates. Rules are transparent, fast, and require no training, but they are brittle
in exactly the dimensions that matter here---dialectal variation, code-switching,
and the combinatorial diversity of natural phrasing. A regex that captures
``paid 50 for groceries'' will miss ``groceries ran me about fifty'' and will be
defeated entirely by Arabic dialect. Rules cannot scale to the linguistic space
of the target users.

\subsection{Supervised Sequence Labeling}

A second approach treats extraction as token classification: label each word of
the transcript with a tag (amount, merchant, date, \dots) using a model such as
BiLSTM-CRF~\cite{lample2016ner} or a fine-tuned transformer encoder like
AraBERT~\cite{antoun2020arabert} or multilingual
BERT/XLM-R~\cite{devlin2019bert,conneau2020xlmr}. This is the classical
named-entity-recognition formulation~\cite{li2022ner} and can be very accurate.
Its decisive drawback for this project is \emph{data}: supervised labeling needs
thousands of annotated, domain-specific, bilingual financial utterances, which do
not exist and would be expensive to create. MiPay therefore reserves this
approach for future work---fine-tuning on user-corrected data, once the product
has accumulated it (Chapter~6).

\subsection{Generative LLM Extraction}

The third approach, and the one adopted, uses a generative instruction-tuned LLM
to read the transcript and emit the structured record directly. With few-shot
prompting~\cite{brown2020gpt3} and constrained
decoding~\cite{willard2023outlines}, this method needs no task-specific training
data, handles dialect and code-switching natively by virtue of the model's
multilingual pre-training, and produces guaranteed-valid structured output. Its
risks---occasional semantic errors and a tendency to perform unreliable
arithmetic such as date math---are contained by the post-processing layer and the
human confirmation step.

\subsection{The Hybrid Approach Adopted by MiPay}
\label{sec:ch1-hybrid}

MiPay's extraction stage is deliberately \emph{hybrid}: a generative LLM does the
linguistically hard, fuzzy work (understanding intent across languages and
mapping it to fields), while a \emph{deterministic Python layer} does the
precise, testable work the LLM does badly. In particular, the LLM is instructed
to copy the \emph{raw} date phrase from the utterance rather than to compute a
calendar date; a deterministic resolver then converts ``yesterday'' or its
dialectal Arabic equivalents into an absolute ISO date anchored on the device's
current date. This division of labor---fuzzy understanding by the model, exact
computation by code---is a recurring theme of the system and a central design
argument of the thesis.

\section{Problems and Challenges}

\begin{description}[leftmargin=1.5em, style=nextline]
    \item[Dialectal and code-switched input.] The pipeline must accept Modern
    Standard Arabic, Gulf/Egyptian/Levantine dialects, English, and arbitrary
    mixtures, where a single sentence may switch language mid-clause.
    \item[Numeral and currency ambiguity.] Amounts may be spoken or written in
    Arabic-Indic numerals, and a single spoken currency word (e.g., \emph{dinar},
    \emph{riyal}) maps to different ISO codes depending on national context.
    \item[Relative-date resolution.] Utterances reference dates relatively
    (``yesterday'', ``last Friday'', dialectal ``\emph{ams}'',
    ``\emph{awwal ams}''); these must be resolved to absolute dates reliably and
    testably---something LLMs do poorly and deterministic code does well.
    \item[Latency on commodity hardware.] The reference deployment is CPU-only
    with 16\,GB of RAM and no GPU, yet the end-to-end voice round-trip must stay
    within an interactive budget.
    \item[Guaranteeing valid output.] The application must never crash on or
    silently corrupt malformed model output, which constrained decoding addresses
    by construction.
    \item[Multi-transaction utterances.] A single sentence may describe several
    transactions (``coffee for 50, a croissant for 100, and my father sent me
    200''), which the pipeline must split into independent records.
    \item[Privacy.] Financial audio is sensitive; the design deletes audio
    immediately after transcription and keeps all inference on the operator's
    machine.
\end{description}

\section{Contributions}
\label{sec:ch1-contributions}

The principal contributions of this thesis are:

\begin{enumerate}[label=\textbf{C\arabic*.}, leftmargin=*]
    \item A complete, reproducible, \textbf{fully self-hosted} bilingual
    voice-to-transaction pipeline assembled from open-source ASR and LLM
    components, with zero per-request cost and no third-party data sharing.
    \item A \textbf{constrained-decoding extraction design} that guarantees
    schema-valid structured output and supports multi-transaction utterances.
    \item A \textbf{deterministic post-processing layer}---numeral normalization,
    dialect-aware relative-date resolution, currency defaulting, and
    confidence-based \texttt{needs\_review} routing---that is independently
    unit-tested and constitutes a reusable, testable contribution.
    \item A \textbf{bilingual evaluation corpus} of more than 260 labeled
    utterances (Arabic, English, and code-switched, including multi-transaction
    cases) and a \textbf{reproducible evaluation harness} reporting Word Error
    Rate and per-field extraction accuracy/F\textsubscript{1} under
    gold-transcript and end-to-end conditions.
    \item A set of \textbf{ablation studies} quantifying the effect of ASR model
    size, LLM size, and few-shot versus zero-shot prompting on extraction
    quality.
    \item A \textbf{cross-platform mobile client} with complete Arabic
    right-to-left and English localization, three input modes, and a
    human-in-the-loop confirmation workflow.
\end{enumerate}

\section{Thesis Organization}

The remainder of this thesis is organized as follows.
\textbf{Chapter~2 (Related Work)} reviews the academic and engineering literature
on speech recognition, large language models, structured extraction and
constrained decoding, and Arabic/multilingual NLP, and positions MiPay relative
to comparable systems.
\textbf{Chapter~3 (System Analysis)} specifies the functional and
non-functional requirements and presents the system models---context, data-flow,
entity-relationship, use-case, class, activity, and sequence diagrams.
\textbf{Chapter~4 (Implementation)} describes the realized system in depth: the
evaluation dataset, the speech-to-text service, the constrained LLM extractor,
the deterministic post-processing layer, the backend and mobile architectures,
and the security model.
\textbf{Chapter~5 (Experimental Results)} defines the evaluation methodology and
metrics and reports transcription and extraction results together with the
ablation studies.
\textbf{Chapter~6 (Conclusion and Future Work)} summarizes the contributions,
states the project's limitations honestly, and lays out a concrete upgrade
path---newer and larger speech and language models, and task-specific
fine-tuning on user-corrected data.
Finally, the \textbf{Appendices} provide a technical implementation reference,
the full extraction prompt and schema, representative source listings, and the
user-interface gallery.
